\begin{frame}[t]{PPO in the LLM Setting}
  Text generation is cast as a sequential decision problem, one token at a time:
  \begin{itemize}\itemsep0.3em
    \item[\textcolor{red}{$\blacktriangleright$}] \textbf{Policy $\pi_\theta$:} the language model itself; $\pi_\theta(a_t \mid s_t)$ is its next-token distribution.
    \item[\textcolor{red}{$\blacktriangleright$}] \textbf{State $s_t$:} the prompt plus every token generated so far.
    \item[\textcolor{red}{$\blacktriangleright$}] \textbf{Action $a_t$:} the next token, drawn from the vocabulary.
    \item[\textcolor{red}{$\blacktriangleright$}] \textbf{Rollout:} one complete sampled response to a prompt.
    \item[\textcolor{red}{$\blacktriangleright$}] \textbf{Reward:} the RM scores the \emph{finished} response, so one number has to be credited back across every token that produced it.
    \item[\textcolor{red}{$\blacktriangleright$}] \textbf{Advantage $\hat{A}_t$:} how much better action $a_t$ turned out than expected at $s_t$. A separate value model, trained alongside the policy, supplies that expectation.
  \end{itemize}
\end{frame}

\begin{frame}[t]{RLHF Training: Proximal Policy Optimization (PPO)}
  \begin{itemize}
    \item[\textcolor{red}{$\blacktriangleright$}] \textbf{Objective:} Maximize the reward given by the reward model (RM).
    \item[\textcolor{red}{$\blacktriangleright$}] \textbf{Reward with KL penalty:} the RM score is shaped so the policy cannot drift far from the SFT model:
      \[
        R(x,y) = r_\phi(x,y) - \beta \log \frac{\pi_\theta(y \mid x)}{\pi_{\text{SFT}}(y \mid x)}
      \]
    \item[\textcolor{red}{$\blacktriangleright$}] \textbf{Clipped surrogate objective}, which PPO \emph{maximises}:
      \[
        L_{\mathrm{PPO}} = \hat{\mathbb{E}}_t\left[
            \min\Big(\rho_t(\theta)\hat{A}_t,\,\,\,\, \text{clip}\big(\rho_t(\theta), 1 - \epsilon, 1 + \epsilon\big)\hat{A}_t\Big)
          \right]
      \]
      where $\rho_t(\theta) = \frac{\pi_\theta(a_t|s_t)}{\pi_\text{old}(a_t|s_t)}$ is the probability ratio (not a reward) and $\hat{A}_t$ is the advantage estimate.
    \item[\textcolor{red}{$\blacktriangleright$}] Clipping caps how far a single update can move the policy, which is what keeps the optimization stable.
  \end{itemize}
\end{frame}

\begin{frame}[t]{PPO --- Pros and Cons}
  \begin{itemize}\itemsep0.5em
    \item[\textcolor{red}{\checkmark}] Stable optimization: the clip bounds every policy update.
    \item[\textcolor{red}{\checkmark}] Learns from the model's own samples, so it can improve on behaviour never demonstrated in the SFT data.
    \item[\textcolor{gray}{\texttimes}] Expensive: each step needs fresh rollouts, and four networks are in play at once, the policy, a frozen SFT copy for the KL term, the reward model, and the value model.
    \item[\textcolor{gray}{\texttimes}] Sensitive to reward model errors: the policy optimizes the RM, not the humans behind it.
    \item[\textcolor{gray}{\texttimes}] Many interacting hyperparameters ($\beta$, $\epsilon$, rollout size, value-loss weight).
  \end{itemize}
\end{frame}
